\section{Conclusion}
We studied whether frozen-VLM failures can guide representation editing without group annotations. \textsc{LatentRisk-VLM} uses task-label margins, class residualization, risk-direction estimation, and candidate-wise inference without updating either encoder. It discovers error-related latent risk directions, not sensitive attributes or demographic groups. Rank-1 editing substantially improves CLIP Waterbirds WGA, while rank-4 editing with worst-class validation selection improves both CLIP Waterbirds and CLIP CelebA over task-accuracy selection. The SigLIP extension and the MetaShift-style cat/dog shift add a more cautious picture: changing the frozen backbone or shift can make task-accuracy selection fail badly, and rank-4 selection may recover WGA only with substantial average-accuracy tradeoffs or modest gains. Post-hoc diagnostics explain why rank-1 CelebA fails: the learned risk direction is stable but poorly aligned with the Male-defined spurious direction, and top-risk examples are not minority-enriched. The revised protocol positions CnC-Frozen, DFR-Frozen, INLP, and LEACE as supervised or attribute-budgeted frozen-feature adaptations rather than as full-backbone training costs. Methods with explicit attribute labels or group-aware validation remain stronger in settings where that supervision is available. Failure-derived directions are useful weak supervision when failure and hidden groups align, not a general substitute for group information.
